A CRC code can be used to detect communication errors. 
It is a cyclic code, and hence generated by a polynomial over a finite field. 
The message is encoded as a string, which is then thought of as a polynomial, 
called the information polynomial. 
Assume that the check polynomial has degree $d$.
The information polynomial is then divided by the check polynomial. 
The remainder is added to the information polynomial multiplied by $X^d.$
This is the codeword, which is sent.

\bigskip

Table~\ref{tab:CRC:process} 
lists the commands for defining a  CRC-process. 
\orbitertableofcommands{crc_process_description}


\bigskip

Table~\ref{tab:CRC:options} 
summarizes options associated with commands for CRC-codes. 
\orbitertableofcommands{crc_options}


\bigskip

Table~\ref{tab:CRC:description} 
lists the available commands for defining a CRC-code. 
\orbitertableofcommands{crc_code_description}


\bigskip


The following command performs an exhaustive search 
over all binary CRC polynomials of degree $k=10$ 
which can detect every error pattern 
of Hamming weight at most $t=3$ in messages of length $n=128.$ 

\sourcecode{CRC_3_128_10}

The program finds 244 polynomials in about 1 minute.



\bigskip





Here is a collection of CRC polynomials from various sources:

\sourcecodevariable{CRC4}


\sourcecodevariable{CRC7}


\sourcecodevariable{CRC8_ATM}


\sourcecodevariable{CRC16_CCITT}


\sourcecodevariable{CRC32_ETHERNET}


\sourcecodevariable{CRC32_CASTAGNOLI}

\sourcecodevariable{CRC64_ECMA182}


\sourcecodevariable{CRC64_ROCKSOFT}



\bigskip


%We test an incorrect implementation of CRC32:
%{\tt
%\input GRAPHICS/THE_MAKEFILE/EXTRACTS/crc32_test.tex
%}
%{\tt
%\input GRAPHICS/THE_MAKEFILE/EXTRACTS/crc32_test_hexdata.tex
%}

We test whether the polynomial crc32 is irreducible:

\sourcecode{crc32_Berlekamp_matrix}


\bigskip

Now, we create some new CRC polynomials over the field $\bbF_{256}.$
To begin with, we create the 771st roots over $\bbF_{256}$:

\sourcecode{CRC_F256_roots_771}

We create a BCH code of length 771 over $\bbF_{256}$ with designed distance 2:

\sourcecode{CRC_F256_BCH_code_d2}

The polynomial in dense coding

\sourcecodevariable{CRC_POLY_Q256_DEG2_DENSE}

We generate C++ source code for the use of this polynomial:

\sourcecode{CRC_F256_BCH_write_code_for_division_d2}

We create a BCH code of length 771 over $\bbF_{256}$ with designed distance 16:

\sourcecode{F256_BCH_code_d16}

The polynomial in sparse coding is:

\sourcecodevariable{POLY_Q256_DEG30_SPARSE}

The polynomial in dense coding is:

\sourcecodevariable{POLY_Q256_DEG30_DENSE}

We generate C++ source code for the use of this polynomial:

\sourcecode{F256_BCH_write_code_for_division_d16}



%The generator polynomial is printed in two ways, sparse and dense. 
%The notion of sparse and dense agrees with that of Section~\ref{sec:vector:builder}.
%Dense means that the coefficient vector of the polynomial is listed in full.
%Sparse means that only the nonzero terms are listed as pairs, the nonzero coefficient and the index of %the term. 
%The coefficient vector determines the generator polynomial 
%$g(X) \in \bbF_{256}[X]$ of the BCH code of length $771$ over $\bbF_{256}.$
%Next, we test if $g(x)$ divides $X^{771}-1$, as it should:
%{\tt
%\input GRAPHICS/THE_MAKEFILE/EXTRACTS/POLY_Q256_DEG30_DENSE.tex
%}
We confirm that the polynomial divides $X^{771}-1$ as it should:

\sourcecode{F256_BCH_code_d16_division}

%This confirms that the remainder after dividing $X^{771}-1$ by $g(X)$ is indeed zero.
The next example introduces three errors. 
The remainder is not zero, so the errors are detected:

\sourcecode{F256_BCH_code_d16_error}

\bigskip

We create the field $\bbF_{256}$:

\sourcecode{F_256}


\bigskip


\sourcecode{CRC_F_16}

\bigskip

\sourcecode{nth_roots_256}

\bigskip

\sourcecode{CRC_F16_roots_51}

\bigskip

\sourcecode{CRC_F16_BCH_code_d3}

\bigskip

\sourcecode{CRC_F16_BCH_code_d3_mindist_dual}

\bigskip

\sourcecode{CRC_F16_BCH_code_d3_dual_weight_enumerator}

\bigskip

\sourcecode{CRC_F16_BCH_code_d3_projective_set}

\bigskip

\sourcecode{CRC_F16_BCH_code_d3_macwilliams}

\bigskip

\sourcecode{CRC_F16_BCH_write_code_for_division_Delta}

\bigskip

\sourcecode{crc_encode_16}

\bigskip

\sourcecode{crc_encode_32}

\bigskip









